\section{Conclusion}\label{sec:conclusion}

We have presented the first comprehensive analysis of network motifs in neural network attribution circuits, establishing that feedforward loop motifs serve as universal structural primitives. Analyzing 200 Neuronpedia attribution graphs across 8 capability domains, we demonstrated FFL universality (median $Z=47.2$, 8/8 domains), functional importance of FFL hubs (1.92$\times$ ablation impact, dose--response $r=0.88$), the superiority of weighted over binary motif features for domain clustering ($NMI=0.705$ vs.\ $0.101$), genuine unique information in motif features ($R^2=0.0184$, MI $retained=23.6\%$), and complete FFL-derivativeness of all 4 universal 4-node motifs. These results bridge network science and mechanistic interpretability, demonstrating that the same structural primitives found in biological regulatory networks \cite{alon2007network, milo2002network} also characterize artificial neural circuits. The practical implications are immediate: FFL motifs and their hub nodes provide structured entry points for circuit discovery, compact fingerprints for cross-model comparison, and principled targets for architecture-aware pruning. The weighted motif features we developed --- particularly FFL path dominance and Onnela intensity --- offer interpretable descriptors that capture both topological and quantitative circuit properties (Table~\ref{table:clustering}; Table~\ref{table:ablation}). Future work should extend this framework to additional models, attribution methods, and dynamic analysis of circuit formation during training. We release all data, code, and analysis scripts for full reproducibility.